% Table tab:mutation_prevalence — TCGA-LUAD analysis cohort mutation
% prevalence (458 tumour slides, 444 patients, 33 sites).
% Counts are SLIDE-level over the 458-slide cohort, verified against the
% run outputs: the KRAS co-activation pass reports 126 KRAS-positive
% slides, while the patient-level OOF adjudication reports 124
% KRAS-positive patients (33+42+49 across the three co-activation
% strata). The superseded 687-image inventory gave systematically lower
% rates (TP53 38.0%, KRAS 20.4%) because it mixed in squamous cases and
% solid-tissue-normal sections.
\small
\setlength{\tabcolsep}{8pt}
\begin{tabular}{l c c c}
\toprule
\textbf{Gene} & \textbf{Positive slides} ($n_{+}$) &
  \textbf{Prevalence} (\%) & \textbf{Conclusive verdict?} \\
\midrule
\textit{TP53}  & 227 & 49.6 & No \\
\textit{KRAS}  & 126 & 27.5 & No \\
\textit{KEAP1} &  85 & 18.6 & No \\
\textit{EGFR}  &  71 & 15.5 & No \\
\textit{STK11} &  59 & 12.9 & No \\
\textit{RBM10} &  32 &  7.0 & No \\
\bottomrule
\end{tabular}
\begin{flushleft}\footnotesize
$n_{+}\!=\!$~number of slides whose patient carries at least one
non-silent somatic mutation called for the gene under the TCGA
harmonised pipeline, over the $458$-slide analysis cohort.
Nine patients contribute more than one section, so the patient-level
counts are marginally lower (\textit{KRAS}, $124$ of $444$); every
downstream statistic is aggregated to the patient before resampling,
and the estimation floor is quoted at the patient level (Methods).
The coincidence in the \textit{RBM10} row is worth disarming: its
\emph{slide} count, $32$, happens to equal the numerical value of the
floor, but the comparison that matters is against its \emph{patient}
count of $30$, which falls below it. \textit{RBM10} is an
estimation-floor case.
These rates are higher than those of the $687$-image inventory
($\textit{TP53}\ 38.0\%$, $\textit{KRAS}\ 20.4\%$) because that
inventory includes squamous-carcinoma cases and solid-tissue-normal
sections, both of which dilute every LUAD mutation rate.
\emph{Conclusive verdict} is the cohort-level signal-presence flag
(mean AUROC $\geq\!0.70$; see Methods, Cross-validation and
statistical analysis). Under nested cross-validation on this cohort no
gene clears it---the study maximum is \textit{TP53} at $0.661$
(\cref{tab:mutation_luad_resultssec})---so the column is uniformly
negative and prevalence does not order the verdict here.
\end{flushleft}
